% Options for packages loaded elsewhere
% Options for packages loaded elsewhere
\PassOptionsToPackage{unicode}{hyperref}
\PassOptionsToPackage{hyphens}{url}
\PassOptionsToPackage{dvipsnames,svgnames,x11names}{xcolor}
%
\documentclass[
  letterpaper,
  DIV=11,
  numbers=noendperiod]{scrartcl}
\usepackage{xcolor}
\usepackage{amsmath,amssymb}
\setcounter{secnumdepth}{-\maxdimen} % remove section numbering
\usepackage{iftex}
\ifPDFTeX
  \usepackage[T1]{fontenc}
  \usepackage[utf8]{inputenc}
  \usepackage{textcomp} % provide euro and other symbols
\else % if luatex or xetex
  \usepackage{unicode-math} % this also loads fontspec
  \defaultfontfeatures{Scale=MatchLowercase}
  \defaultfontfeatures[\rmfamily]{Ligatures=TeX,Scale=1}
\fi
\usepackage{lmodern}
\ifPDFTeX\else
  % xetex/luatex font selection
\fi
% Use upquote if available, for straight quotes in verbatim environments
\IfFileExists{upquote.sty}{\usepackage{upquote}}{}
\IfFileExists{microtype.sty}{% use microtype if available
  \usepackage[]{microtype}
  \UseMicrotypeSet[protrusion]{basicmath} % disable protrusion for tt fonts
}{}
\makeatletter
\@ifundefined{KOMAClassName}{% if non-KOMA class
  \IfFileExists{parskip.sty}{%
    \usepackage{parskip}
  }{% else
    \setlength{\parindent}{0pt}
    \setlength{\parskip}{6pt plus 2pt minus 1pt}}
}{% if KOMA class
  \KOMAoptions{parskip=half}}
\makeatother
% Make \paragraph and \subparagraph free-standing
\makeatletter
\ifx\paragraph\undefined\else
  \let\oldparagraph\paragraph
  \renewcommand{\paragraph}{
    \@ifstar
      \xxxParagraphStar
      \xxxParagraphNoStar
  }
  \newcommand{\xxxParagraphStar}[1]{\oldparagraph*{#1}\mbox{}}
  \newcommand{\xxxParagraphNoStar}[1]{\oldparagraph{#1}\mbox{}}
\fi
\ifx\subparagraph\undefined\else
  \let\oldsubparagraph\subparagraph
  \renewcommand{\subparagraph}{
    \@ifstar
      \xxxSubParagraphStar
      \xxxSubParagraphNoStar
  }
  \newcommand{\xxxSubParagraphStar}[1]{\oldsubparagraph*{#1}\mbox{}}
  \newcommand{\xxxSubParagraphNoStar}[1]{\oldsubparagraph{#1}\mbox{}}
\fi
\makeatother

\usepackage{color}
\usepackage{fancyvrb}
\newcommand{\VerbBar}{|}
\newcommand{\VERB}{\Verb[commandchars=\\\{\}]}
\DefineVerbatimEnvironment{Highlighting}{Verbatim}{commandchars=\\\{\}}
% Add ',fontsize=\small' for more characters per line
\usepackage{framed}
\definecolor{shadecolor}{RGB}{241,243,245}
\newenvironment{Shaded}{\begin{snugshade}}{\end{snugshade}}
\newcommand{\AlertTok}[1]{\textcolor[rgb]{0.68,0.00,0.00}{#1}}
\newcommand{\AnnotationTok}[1]{\textcolor[rgb]{0.37,0.37,0.37}{#1}}
\newcommand{\AttributeTok}[1]{\textcolor[rgb]{0.40,0.45,0.13}{#1}}
\newcommand{\BaseNTok}[1]{\textcolor[rgb]{0.68,0.00,0.00}{#1}}
\newcommand{\BuiltInTok}[1]{\textcolor[rgb]{0.00,0.23,0.31}{#1}}
\newcommand{\CharTok}[1]{\textcolor[rgb]{0.13,0.47,0.30}{#1}}
\newcommand{\CommentTok}[1]{\textcolor[rgb]{0.37,0.37,0.37}{#1}}
\newcommand{\CommentVarTok}[1]{\textcolor[rgb]{0.37,0.37,0.37}{\textit{#1}}}
\newcommand{\ConstantTok}[1]{\textcolor[rgb]{0.56,0.35,0.01}{#1}}
\newcommand{\ControlFlowTok}[1]{\textcolor[rgb]{0.00,0.23,0.31}{\textbf{#1}}}
\newcommand{\DataTypeTok}[1]{\textcolor[rgb]{0.68,0.00,0.00}{#1}}
\newcommand{\DecValTok}[1]{\textcolor[rgb]{0.68,0.00,0.00}{#1}}
\newcommand{\DocumentationTok}[1]{\textcolor[rgb]{0.37,0.37,0.37}{\textit{#1}}}
\newcommand{\ErrorTok}[1]{\textcolor[rgb]{0.68,0.00,0.00}{#1}}
\newcommand{\ExtensionTok}[1]{\textcolor[rgb]{0.00,0.23,0.31}{#1}}
\newcommand{\FloatTok}[1]{\textcolor[rgb]{0.68,0.00,0.00}{#1}}
\newcommand{\FunctionTok}[1]{\textcolor[rgb]{0.28,0.35,0.67}{#1}}
\newcommand{\ImportTok}[1]{\textcolor[rgb]{0.00,0.46,0.62}{#1}}
\newcommand{\InformationTok}[1]{\textcolor[rgb]{0.37,0.37,0.37}{#1}}
\newcommand{\KeywordTok}[1]{\textcolor[rgb]{0.00,0.23,0.31}{\textbf{#1}}}
\newcommand{\NormalTok}[1]{\textcolor[rgb]{0.00,0.23,0.31}{#1}}
\newcommand{\OperatorTok}[1]{\textcolor[rgb]{0.37,0.37,0.37}{#1}}
\newcommand{\OtherTok}[1]{\textcolor[rgb]{0.00,0.23,0.31}{#1}}
\newcommand{\PreprocessorTok}[1]{\textcolor[rgb]{0.68,0.00,0.00}{#1}}
\newcommand{\RegionMarkerTok}[1]{\textcolor[rgb]{0.00,0.23,0.31}{#1}}
\newcommand{\SpecialCharTok}[1]{\textcolor[rgb]{0.37,0.37,0.37}{#1}}
\newcommand{\SpecialStringTok}[1]{\textcolor[rgb]{0.13,0.47,0.30}{#1}}
\newcommand{\StringTok}[1]{\textcolor[rgb]{0.13,0.47,0.30}{#1}}
\newcommand{\VariableTok}[1]{\textcolor[rgb]{0.07,0.07,0.07}{#1}}
\newcommand{\VerbatimStringTok}[1]{\textcolor[rgb]{0.13,0.47,0.30}{#1}}
\newcommand{\WarningTok}[1]{\textcolor[rgb]{0.37,0.37,0.37}{\textit{#1}}}

\usepackage{longtable,booktabs,array}
\newcounter{none} % for unnumbered tables
\usepackage{calc} % for calculating minipage widths
% Correct order of tables after \paragraph or \subparagraph
\usepackage{etoolbox}
\makeatletter
\patchcmd\longtable{\par}{\if@noskipsec\mbox{}\fi\par}{}{}
\makeatother
% Allow footnotes in longtable head/foot
\IfFileExists{footnotehyper.sty}{\usepackage{footnotehyper}}{\usepackage{footnote}}
\makesavenoteenv{longtable}
\usepackage{graphicx}
\makeatletter
\newsavebox\pandoc@box
\newcommand*\pandocbounded[1]{% scales image to fit in text height/width
  \sbox\pandoc@box{#1}%
  \Gscale@div\@tempa{\textheight}{\dimexpr\ht\pandoc@box+\dp\pandoc@box\relax}%
  \Gscale@div\@tempb{\linewidth}{\wd\pandoc@box}%
  \ifdim\@tempb\p@<\@tempa\p@\let\@tempa\@tempb\fi% select the smaller of both
  \ifdim\@tempa\p@<\p@\scalebox{\@tempa}{\usebox\pandoc@box}%
  \else\usebox{\pandoc@box}%
  \fi%
}
% Set default figure placement to htbp
\def\fps@figure{htbp}
\makeatother





\setlength{\emergencystretch}{3em} % prevent overfull lines

\providecommand{\tightlist}{%
  \setlength{\itemsep}{0pt}\setlength{\parskip}{0pt}}



 


\KOMAoption{captions}{tableheading}
\makeatletter
\@ifpackageloaded{caption}{}{\usepackage{caption}}
\AtBeginDocument{%
\ifdefined\contentsname
  \renewcommand*\contentsname{Table of contents}
\else
  \newcommand\contentsname{Table of contents}
\fi
\ifdefined\listfigurename
  \renewcommand*\listfigurename{List of Figures}
\else
  \newcommand\listfigurename{List of Figures}
\fi
\ifdefined\listtablename
  \renewcommand*\listtablename{List of Tables}
\else
  \newcommand\listtablename{List of Tables}
\fi
\ifdefined\figurename
  \renewcommand*\figurename{Figure}
\else
  \newcommand\figurename{Figure}
\fi
\ifdefined\tablename
  \renewcommand*\tablename{Table}
\else
  \newcommand\tablename{Table}
\fi
}
\@ifpackageloaded{float}{}{\usepackage{float}}
\floatstyle{ruled}
\@ifundefined{c@chapter}{\newfloat{codelisting}{h}{lop}}{\newfloat{codelisting}{h}{lop}[chapter]}
\floatname{codelisting}{Listing}
\newcommand*\listoflistings{\listof{codelisting}{List of Listings}}
\makeatother
\makeatletter
\makeatother
\makeatletter
\@ifpackageloaded{caption}{}{\usepackage{caption}}
\@ifpackageloaded{subcaption}{}{\usepackage{subcaption}}
\makeatother
\usepackage{bookmark}
\IfFileExists{xurl.sty}{\usepackage{xurl}}{} % add URL line breaks if available
\urlstyle{same}
% fallback for those not using the hyperref driver hyperxmp:
\makeatletter
\@ifundefined{xmpquote}{\newcommand{\xmpquote}[1]{#1}}{}
\makeatother
\hypersetup{
  pdftitle={hw1},
  pdfauthor={TING-YU CHANG},
  colorlinks=true,
  linkcolor={blue},
  filecolor={Maroon},
  citecolor={Blue},
  urlcolor={Blue},
  pdfcreator={LaTeX via pandoc}}


\title{hw1}
\author{TING-YU CHANG}
\date{}
\begin{document}
\maketitle


This is the link to my repo on
github:https://github.com/tingyuuucc/stats506-HW1

\subsection{Q1:Exploratory Analysis of Soybean
Cultivars}\label{q1exploratory-analysis-of-soybean-cultivars}

\subsubsection{a. Import the data into a data.frame in R. Use the
dataset documentation to assign appropriate variable names. (Note:
Downloading and unzipping the file can take place outside of your
submitted document, but importing the file should be in the
submission.)}\label{a.-import-the-data-into-a-data.frame-in-r.-use-the-dataset-documentation-to-assign-appropriate-variable-names.-note-downloading-and-unzipping-the-file-can-take-place-outside-of-your-submitted-document-but-importing-the-file-should-be-in-the-submission.}

\begin{Shaded}
\begin{Highlighting}[]
\NormalTok{soybean }\OtherTok{\textless{}{-}} \FunctionTok{read.csv}\NormalTok{(}\StringTok{"data.csv"}\NormalTok{)}

\FunctionTok{colnames}\NormalTok{(soybean) }\OtherTok{\textless{}{-}} \FunctionTok{c}\NormalTok{(}
  \StringTok{"Season"}\NormalTok{,}
  \StringTok{"Cultivar"}\NormalTok{,}
  \StringTok{"Repetition"}\NormalTok{,}
  \StringTok{"PH"}\NormalTok{,}
  \StringTok{"IFP"}\NormalTok{,}
  \StringTok{"NLP"}\NormalTok{,}
  \StringTok{"NGP"}\NormalTok{,}
  \StringTok{"NGL"}\NormalTok{,}
  \StringTok{"NS"}\NormalTok{,}
  \StringTok{"MHG"}\NormalTok{,}
  \StringTok{"GY"}
\NormalTok{)}
\end{Highlighting}
\end{Shaded}

\subsubsection{b. Using the information provided in the dataset
documentation: Verify that there are 320 observations in the dataset; 40
distinct cultivars; and all the observations were collected from two
seasons. Produce a table showing the number of observations from each
season.}\label{b.-using-the-information-provided-in-the-dataset-documentation-verify-that-there-are-320-observations-in-the-dataset-40-distinct-cultivars-and-all-the-observations-were-collected-from-two-seasons.-produce-a-table-showing-the-number-of-observations-from-each-season.}

\begin{Shaded}
\begin{Highlighting}[]
\CommentTok{\# number of obs.}
\FunctionTok{nrow}\NormalTok{(soybean)}
\end{Highlighting}
\end{Shaded}

\begin{verbatim}
[1] 320
\end{verbatim}

\begin{Shaded}
\begin{Highlighting}[]
\CommentTok{\# number of distinct cultivars}
\FunctionTok{length}\NormalTok{(}\FunctionTok{unique}\NormalTok{(soybean}\SpecialCharTok{$}\NormalTok{Cultivar))}
\end{Highlighting}
\end{Shaded}

\begin{verbatim}
[1] 40
\end{verbatim}

\begin{Shaded}
\begin{Highlighting}[]
\CommentTok{\# number of distinct seasons}
\FunctionTok{length}\NormalTok{(}\FunctionTok{unique}\NormalTok{(soybean}\SpecialCharTok{$}\NormalTok{Season))}
\end{Highlighting}
\end{Shaded}

\begin{verbatim}
[1] 2
\end{verbatim}

\begin{Shaded}
\begin{Highlighting}[]
\CommentTok{\# number of observations from each season}
\FunctionTok{table}\NormalTok{(soybean}\SpecialCharTok{$}\NormalTok{Season)}
\end{Highlighting}
\end{Shaded}

\begin{verbatim}

  1   2 
160 160 
\end{verbatim}

The dataset contains 320 observations, 40 distinct cultivars, and two
seasons. Each season contains 160 observations.

\subsubsection{c.~Use the data to answer the following
questions}\label{c.-use-the-data-to-answer-the-following-questions}

\subsubsection{1.What is the correlation between plant height (PH) and
grain yield
(GY)?}\label{what-is-the-correlation-between-plant-height-ph-and-grain-yield-gy}

\begin{Shaded}
\begin{Highlighting}[]
\FunctionTok{cor}\NormalTok{(soybean}\SpecialCharTok{$}\NormalTok{PH, soybean}\SpecialCharTok{$}\NormalTok{GY)}
\end{Highlighting}
\end{Shaded}

\begin{verbatim}
[1] 0.1232807
\end{verbatim}

The correlation between plant height (PH) and grain yield (GY) is
0.1232807.

\subsubsection{2.Calculate the correlation between PH and GY separately
for each season. Which season has the higher
correlation?}\label{calculate-the-correlation-between-ph-and-gy-separately-for-each-season.-which-season-has-the-higher-correlation}

\begin{Shaded}
\begin{Highlighting}[]
\FunctionTok{cor}\NormalTok{(soybean}\SpecialCharTok{$}\NormalTok{PH[soybean}\SpecialCharTok{$}\NormalTok{Season}\SpecialCharTok{==}\DecValTok{1}\NormalTok{],}
\NormalTok{    soybean}\SpecialCharTok{$}\NormalTok{GY[soybean}\SpecialCharTok{$}\NormalTok{Season}\SpecialCharTok{==}\DecValTok{1}\NormalTok{])}
\end{Highlighting}
\end{Shaded}

\begin{verbatim}
[1] 0.2308581
\end{verbatim}

\begin{Shaded}
\begin{Highlighting}[]
\FunctionTok{cor}\NormalTok{(soybean}\SpecialCharTok{$}\NormalTok{PH[soybean}\SpecialCharTok{$}\NormalTok{Season}\SpecialCharTok{==}\DecValTok{2}\NormalTok{],}
\NormalTok{    soybean}\SpecialCharTok{$}\NormalTok{GY[soybean}\SpecialCharTok{$}\NormalTok{Season}\SpecialCharTok{==}\DecValTok{2}\NormalTok{])}
\end{Highlighting}
\end{Shaded}

\begin{verbatim}
[1] -0.0227046
\end{verbatim}

Season 2 has the higher correlation.

\subsubsection{3.What is the grain yield of the soybean observation with
the greatest plant
height?}\label{what-is-the-grain-yield-of-the-soybean-observation-with-the-greatest-plant-height}

\begin{Shaded}
\begin{Highlighting}[]
\NormalTok{soybean[soybean}\SpecialCharTok{$}\NormalTok{PH }\SpecialCharTok{==} \FunctionTok{max}\NormalTok{(soybean}\SpecialCharTok{$}\NormalTok{PH), }\FunctionTok{c}\NormalTok{(}\StringTok{"PH"}\NormalTok{,}\StringTok{"GY"}\NormalTok{)]}
\end{Highlighting}
\end{Shaded}

\begin{verbatim}
      PH       GY
315 94.8 3304.678
\end{verbatim}

The grain yield of the soybean observation with the greatest plant
height = 3304.678.

\subsubsection{4.What percentage of observations have a thousand seed
weight (MHG) greater than the overall average thousand seed
weight?}\label{what-percentage-of-observations-have-a-thousand-seed-weight-mhg-greater-than-the-overall-average-thousand-seed-weight}

\begin{Shaded}
\begin{Highlighting}[]
\NormalTok{mean\_MHG }\OtherTok{\textless{}{-}} \FunctionTok{mean}\NormalTok{(soybean}\SpecialCharTok{$}\NormalTok{MHG)}
\CommentTok{\# covert to \%}
\FunctionTok{mean}\NormalTok{(soybean}\SpecialCharTok{$}\NormalTok{MHG }\SpecialCharTok{\textgreater{}}\NormalTok{ mean\_MHG) }\SpecialCharTok{*} \DecValTok{100}
\end{Highlighting}
\end{Shaded}

\begin{verbatim}
[1] 47.1875
\end{verbatim}

47.1875\%

\subsubsection{5.Which cultivar has the highest average grain
yield?}\label{which-cultivar-has-the-highest-average-grain-yield}

\begin{Shaded}
\begin{Highlighting}[]
\CommentTok{\# calculate average GY}
\NormalTok{avg\_GY }\OtherTok{\textless{}{-}} \FunctionTok{aggregate}\NormalTok{(GY }\SpecialCharTok{\textasciitilde{}}\NormalTok{ Cultivar,}
                    \AttributeTok{data =}\NormalTok{ soybean,}
                    \AttributeTok{FUN =}\NormalTok{ mean)}
\CommentTok{\# find the highest average GY}
\NormalTok{avg\_GY[}\FunctionTok{which.max}\NormalTok{(avg\_GY}\SpecialCharTok{$}\NormalTok{GY), ]}
\end{Highlighting}
\end{Shaded}

\begin{verbatim}
  Cultivar       GY
9 98R30 CE 3964.968
\end{verbatim}

The answer is 98R30 CE.

\subsubsection{d.~Create a table showing the average value of each
quantitative variable: One row containing overall averages. One row for
each season containing season-specific averages. The table should
clearly identify what each row represents. (This table does not need to
be ``fancy'' but should clearly identify what each value
represents.)}\label{d.-create-a-table-showing-the-average-value-of-each-quantitative-variable-one-row-containing-overall-averages.-one-row-for-each-season-containing-season-specific-averages.-the-table-should-clearly-identify-what-each-row-represents.-this-table-does-not-need-to-be-fancy-but-should-clearly-identify-what-each-value-represents.}

\begin{Shaded}
\begin{Highlighting}[]
\CommentTok{\# quantitative variables}
\NormalTok{quant\_vars }\OtherTok{\textless{}{-}} \FunctionTok{c}\NormalTok{(}\StringTok{"PH"}\NormalTok{,}\StringTok{"IFP"}\NormalTok{,}\StringTok{"NGP"}\NormalTok{,}\StringTok{"NGL"}\NormalTok{,}\StringTok{"NS"}\NormalTok{,}\StringTok{"MHG"}\NormalTok{,}\StringTok{"GY"}\NormalTok{)}

\CommentTok{\# overall averages}
\NormalTok{overall\_avg }\OtherTok{\textless{}{-}} \FunctionTok{data.frame}\NormalTok{(}
  \AttributeTok{Group =} \StringTok{"Overall"}\NormalTok{,}
  \FunctionTok{t}\NormalTok{(}\FunctionTok{colMeans}\NormalTok{(soybean[quant\_vars]))}
\NormalTok{)}

\CommentTok{\# season{-}specific averages}
\NormalTok{season\_avg }\OtherTok{\textless{}{-}} \FunctionTok{aggregate}\NormalTok{(}
\NormalTok{  soybean[quant\_vars],}
  \AttributeTok{by =} \FunctionTok{list}\NormalTok{(}\AttributeTok{Group =} \FunctionTok{paste0}\NormalTok{(}\StringTok{"Season "}\NormalTok{, soybean}\SpecialCharTok{$}\NormalTok{Season)),}
  \AttributeTok{FUN =}\NormalTok{ mean}
\NormalTok{)}

\CommentTok{\# ChatGPT helped me figure out a more efficient way to code my loop for problem d by doing aggregate.}

\CommentTok{\#combine into one table}
\NormalTok{avg\_table }\OtherTok{\textless{}{-}} \FunctionTok{rbind}\NormalTok{(overall\_avg, season\_avg)}

\NormalTok{avg\_table}
\end{Highlighting}
\end{Shaded}

\begin{verbatim}
     Group       PH    IFP      NGP      NGL       NS      MHG       GY
1  Overall 68.38678 15.465 135.0858 2.290844 4.071656 168.3223 3418.554
2 Season 1 65.98506 16.485 135.2350 2.299688 4.816625 162.2888 3428.452
3 Season 2 70.78850 14.445 134.9367 2.282000 3.326687 174.3558 3408.655
\end{verbatim}

\subsubsection{e. Carry out a series of t-tests to examine whether the
ngrain yield (GY) differs between the two seasons. State the null and
the alternative hypotheses. Present the R output and interpret the
results.}\label{e.-carry-out-a-series-of-t-tests-to-examine-whether-the-ngrain-yield-gy-differs-between-the-two-seasons.-state-the-null-and-the-alternative-hypotheses.-present-the-r-output-and-interpret-the-results.}

\begin{Shaded}
\begin{Highlighting}[]
\CommentTok{\# calculate GY by season}
\NormalTok{GY1 }\OtherTok{\textless{}{-}}\NormalTok{ soybean}\SpecialCharTok{$}\NormalTok{GY[soybean}\SpecialCharTok{$}\NormalTok{Season }\SpecialCharTok{==} \DecValTok{1}\NormalTok{]}
\NormalTok{GY2 }\OtherTok{\textless{}{-}}\NormalTok{ soybean}\SpecialCharTok{$}\NormalTok{GY[soybean}\SpecialCharTok{$}\NormalTok{Season }\SpecialCharTok{==} \DecValTok{2}\NormalTok{]}

\CommentTok{\# sample statistics}
\NormalTok{n1 }\OtherTok{\textless{}{-}} \FunctionTok{length}\NormalTok{(GY1)}
\NormalTok{n2 }\OtherTok{\textless{}{-}} \FunctionTok{length}\NormalTok{(GY2)}

\NormalTok{mean1 }\OtherTok{\textless{}{-}} \FunctionTok{mean}\NormalTok{(GY1)}
\NormalTok{mean2 }\OtherTok{\textless{}{-}} \FunctionTok{mean}\NormalTok{(GY2)}

\NormalTok{var1 }\OtherTok{\textless{}{-}} \FunctionTok{var}\NormalTok{(GY1)}
\NormalTok{var2 }\OtherTok{\textless{}{-}} \FunctionTok{var}\NormalTok{(GY2)}

\CommentTok{\# calculate t statistics}
\NormalTok{t\_stat }\OtherTok{\textless{}{-}}\NormalTok{ (mean1 }\SpecialCharTok{{-}}\NormalTok{ mean2) }\SpecialCharTok{/} \FunctionTok{sqrt}\NormalTok{(var1 }\SpecialCharTok{/}\NormalTok{ n1 }\SpecialCharTok{+}\NormalTok{ var2 }\SpecialCharTok{/}\NormalTok{ n2)}

\CommentTok{\# calculate df}
\NormalTok{df }\OtherTok{\textless{}{-}}\NormalTok{ (var1 }\SpecialCharTok{/}\NormalTok{ n1 }\SpecialCharTok{+}\NormalTok{ var2 }\SpecialCharTok{/}\NormalTok{ n2)}\SpecialCharTok{\^{}}\DecValTok{2} \SpecialCharTok{/}\NormalTok{ ((var1 }\SpecialCharTok{/}\NormalTok{ n1)}\SpecialCharTok{\^{}}\DecValTok{2} \SpecialCharTok{/}\NormalTok{ (n1 }\SpecialCharTok{{-}} \DecValTok{1}\NormalTok{) }\SpecialCharTok{+}\NormalTok{ (var2 }\SpecialCharTok{/}\NormalTok{ n2)}\SpecialCharTok{\^{}}\DecValTok{2} \SpecialCharTok{/}\NormalTok{ (n2 }\SpecialCharTok{{-}} \DecValTok{1}\NormalTok{))}

\CommentTok{\# calculate two{-}sided p{-}value}
\NormalTok{p\_value }\OtherTok{\textless{}{-}} \DecValTok{2} \SpecialCharTok{*} \FunctionTok{pt}\NormalTok{(}\SpecialCharTok{{-}}\FunctionTok{abs}\NormalTok{(t\_stat), }\AttributeTok{df =}\NormalTok{ df)}

\NormalTok{t\_stat}
\end{Highlighting}
\end{Shaded}

\begin{verbatim}
[1] 0.351546
\end{verbatim}

\begin{Shaded}
\begin{Highlighting}[]
\NormalTok{df}
\end{Highlighting}
\end{Shaded}

\begin{verbatim}
[1] 231.8466
\end{verbatim}

\begin{Shaded}
\begin{Highlighting}[]
\NormalTok{p\_value}
\end{Highlighting}
\end{Shaded}

\begin{verbatim}
[1] 0.725498
\end{verbatim}

\begin{Shaded}
\begin{Highlighting}[]
\FunctionTok{t.test}\NormalTok{(GY }\SpecialCharTok{\textasciitilde{}}\NormalTok{ Season, }\AttributeTok{data =}\NormalTok{ soybean)}
\end{Highlighting}
\end{Shaded}

\begin{verbatim}

    Welch Two Sample t-test

data:  GY by Season
t = 0.35155, df = 231.85, p-value = 0.7255
alternative hypothesis: true difference in means between group 1 and group 2 is not equal to 0
95 percent confidence interval:
 -91.15712 130.75170
sample estimates:
mean in group 1 mean in group 2 
       3428.452        3408.655 
\end{verbatim}

Null hypothesis: \(H_0: \mu_1 = \mu_2\), meaning that the mean grain
yield is the same in the two seasons. Alternative hypothesis:
\(H_0: \mu_1 neq \mu_2\), meaning that the mean grain yield differs
between the two seasons. The Welch two-sample t-test gave (t = 0.35155),
(df = 231.85), and a p-value of 0.7255. Since the p-value is greater
than 0.05, we fail to reject the null hypothesis. There is not
sufficient evidence to conclude that the mean grain yield differs
between the two seasons.

\section{ChatGPT helped me figure out a clearer way to show my result of
e by using
latex.}\label{chatgpt-helped-me-figure-out-a-clearer-way-to-show-my-result-of-e-by-using-latex.}

\subsection{Q2:Adult Income Dataset}\label{q2adult-income-dataset}

\subsubsection{a. Import the data into a data.frame in R. As with the
previous data, you may download the data outside of your submission, but
importation should take place inside the problem set submission.
ß}\label{a.-import-the-data-into-a-data.frame-in-r.-as-with-the-previous-data-you-may-download-the-data-outside-of-your-submission-but-importation-should-take-place-inside-the-problem-set-submission.-uxdf}

\begin{Shaded}
\begin{Highlighting}[]
\CommentTok{\# import the dataset}
\NormalTok{adult }\OtherTok{\textless{}{-}} \FunctionTok{read.csv}\NormalTok{(}
  \StringTok{"adult/adult.data"}\NormalTok{,}
  \AttributeTok{header =} \ConstantTok{FALSE}\NormalTok{,}
  \AttributeTok{strip.white =} \ConstantTok{TRUE}\NormalTok{,}
  \AttributeTok{na.string =} \StringTok{"?"}
\NormalTok{)}

\CommentTok{\# assign variable name}
\FunctionTok{colnames}\NormalTok{(adult) }\OtherTok{\textless{}{-}} \FunctionTok{c}\NormalTok{(}
  \StringTok{"age"}\NormalTok{,}
  \StringTok{"workclass"}\NormalTok{,}
  \StringTok{"fnlwgt"}\NormalTok{,}
  \StringTok{"education"}\NormalTok{,}
  \StringTok{"education\_num"}\NormalTok{,}
  \StringTok{"marital\_status"}\NormalTok{,}
  \StringTok{"occupation"}\NormalTok{,}
  \StringTok{"relationship"}\NormalTok{,}
  \StringTok{"race"}\NormalTok{,}
  \StringTok{"sex"}\NormalTok{,}
  \StringTok{"capital\_gain"}\NormalTok{,}
  \StringTok{"capital\_loss"}\NormalTok{,}
  \StringTok{"hours\_per\_week"}\NormalTok{,}
  \StringTok{"native\_country"}\NormalTok{,}
  \StringTok{"income"}
\NormalTok{)}

\CommentTok{\# Preview the first few obs.}
\FunctionTok{head}\NormalTok{(adult)}
\end{Highlighting}
\end{Shaded}

\begin{verbatim}
  age        workclass fnlwgt education education_num     marital_status
1  39        State-gov  77516 Bachelors            13      Never-married
2  50 Self-emp-not-inc  83311 Bachelors            13 Married-civ-spouse
3  38          Private 215646   HS-grad             9           Divorced
4  53          Private 234721      11th             7 Married-civ-spouse
5  28          Private 338409 Bachelors            13 Married-civ-spouse
6  37          Private 284582   Masters            14 Married-civ-spouse
         occupation  relationship  race    sex capital_gain capital_loss
1      Adm-clerical Not-in-family White   Male         2174            0
2   Exec-managerial       Husband White   Male            0            0
3 Handlers-cleaners Not-in-family White   Male            0            0
4 Handlers-cleaners       Husband Black   Male            0            0
5    Prof-specialty          Wife Black Female            0            0
6   Exec-managerial          Wife White Female            0            0
  hours_per_week native_country income
1             40  United-States  <=50K
2             13  United-States  <=50K
3             40  United-States  <=50K
4             40  United-States  <=50K
5             40           Cuba  <=50K
6             40  United-States  <=50K
\end{verbatim}

\begin{Shaded}
\begin{Highlighting}[]
\CommentTok{\# check structure}
\FunctionTok{str}\NormalTok{(adult)}
\end{Highlighting}
\end{Shaded}

\begin{verbatim}
'data.frame':   32561 obs. of  15 variables:
 $ age           : int  39 50 38 53 28 37 49 52 31 42 ...
 $ workclass     : chr  "State-gov" "Self-emp-not-inc" "Private" "Private" ...
 $ fnlwgt        : int  77516 83311 215646 234721 338409 284582 160187 209642 45781 159449 ...
 $ education     : chr  "Bachelors" "Bachelors" "HS-grad" "11th" ...
 $ education_num : int  13 13 9 7 13 14 5 9 14 13 ...
 $ marital_status: chr  "Never-married" "Married-civ-spouse" "Divorced" "Married-civ-spouse" ...
 $ occupation    : chr  "Adm-clerical" "Exec-managerial" "Handlers-cleaners" "Handlers-cleaners" ...
 $ relationship  : chr  "Not-in-family" "Husband" "Not-in-family" "Husband" ...
 $ race          : chr  "White" "White" "White" "Black" ...
 $ sex           : chr  "Male" "Male" "Male" "Male" ...
 $ capital_gain  : int  2174 0 0 0 0 0 0 0 14084 5178 ...
 $ capital_loss  : int  0 0 0 0 0 0 0 0 0 0 ...
 $ hours_per_week: int  40 13 40 40 40 40 16 45 50 40 ...
 $ native_country: chr  "United-States" "United-States" "United-States" "United-States" ...
 $ income        : chr  "<=50K" "<=50K" "<=50K" "<=50K" ...
\end{verbatim}

\subsubsection{b.Clean and simplify the variable
names.}\label{b.clean-and-simplify-the-variable-names.}

\begin{Shaded}
\begin{Highlighting}[]
\FunctionTok{names}\NormalTok{(adult) }\OtherTok{\textless{}{-}} \FunctionTok{c}\NormalTok{(}
  \StringTok{"age"}\NormalTok{,}
  \StringTok{"workclass"}\NormalTok{,}
  \StringTok{"weight"}\NormalTok{,}
  \StringTok{"education"}\NormalTok{,}
  \StringTok{"educ\_num"}\NormalTok{,}
  \StringTok{"marital"}\NormalTok{,}
  \StringTok{"occupation"}\NormalTok{,}
  \StringTok{"relationship"}\NormalTok{,}
  \StringTok{"race"}\NormalTok{,}
  \StringTok{"sex"}\NormalTok{,}
  \StringTok{"cap\_gain"}\NormalTok{,}
  \StringTok{"cap\_loss"}\NormalTok{,}
  \StringTok{"hours"}\NormalTok{,}
  \StringTok{"country"}\NormalTok{,}
  \StringTok{"income"}
\NormalTok{)}

\FunctionTok{names}\NormalTok{(adult)}
\end{Highlighting}
\end{Shaded}

\begin{verbatim}
 [1] "age"          "workclass"    "weight"       "education"    "educ_num"    
 [6] "marital"      "occupation"   "relationship" "race"         "sex"         
[11] "cap_gain"     "cap_loss"     "hours"        "country"      "income"      
\end{verbatim}

\subsubsection{c.Examine missing values. Report the number of missing
observations for each
variable.}\label{c.examine-missing-values.-report-the-number-of-missing-observations-for-each-variable.}

\begin{Shaded}
\begin{Highlighting}[]
\CommentTok{\# create a table of missing observations in each variable}
\NormalTok{missing\_table }\OtherTok{\textless{}{-}} \FunctionTok{data.frame}\NormalTok{(}
  \AttributeTok{Variable =} \FunctionTok{names}\NormalTok{(adult),}
  \AttributeTok{Missing =} \FunctionTok{colSums}\NormalTok{(}\FunctionTok{is.na}\NormalTok{(adult))}
\NormalTok{)}

\CommentTok{\# display the number}
\NormalTok{missing\_table}
\end{Highlighting}
\end{Shaded}

\begin{verbatim}
                 Variable Missing
age                   age       0
workclass       workclass    1836
weight             weight       0
education       education       0
educ_num         educ_num       0
marital           marital       0
occupation     occupation    1843
relationship relationship       0
race                 race       0
sex                   sex       0
cap_gain         cap_gain       0
cap_loss         cap_loss       0
hours               hours       0
country           country     583
income             income       0
\end{verbatim}

Missing values are present in `workclass', `occupation', and `country',
with 1836,1843, and 583 missing observations, respectively.

\subsubsection{d.~Restrict the sample to individuals working at least 10
hours per week. Confirm the sample size before and after
filtering.}\label{d.-restrict-the-sample-to-individuals-working-at-least-10-hours-per-week.-confirm-the-sample-size-before-and-after-filtering.}

\begin{Shaded}
\begin{Highlighting}[]
\CommentTok{\# check the sample size}
\NormalTok{n\_before }\OtherTok{\textless{}{-}} \FunctionTok{nrow}\NormalTok{(adult)}
\NormalTok{n\_before}
\end{Highlighting}
\end{Shaded}

\begin{verbatim}
[1] 32561
\end{verbatim}

\begin{Shaded}
\begin{Highlighting}[]
\CommentTok{\# restrict the sample}
\NormalTok{adult }\OtherTok{\textless{}{-}}\NormalTok{ adult[adult}\SpecialCharTok{$}\NormalTok{hours }\SpecialCharTok{\textgreater{}=} \DecValTok{10}\NormalTok{, ]}

\CommentTok{\# check the sample size after filtering}
\NormalTok{n\_after }\OtherTok{\textless{}{-}} \FunctionTok{nrow}\NormalTok{(adult)}
\NormalTok{n\_after}
\end{Highlighting}
\end{Shaded}

\begin{verbatim}
[1] 32103
\end{verbatim}

\begin{Shaded}
\begin{Highlighting}[]
\CommentTok{\# create the table showing the sample size before and after}
\NormalTok{sample\_size }\OtherTok{\textless{}{-}} \FunctionTok{data.frame}\NormalTok{(}
  \AttributeTok{Stage =} \FunctionTok{c}\NormalTok{(}\StringTok{"Before"}\NormalTok{, }\StringTok{"After filtering"}\NormalTok{),}
  \AttributeTok{Sample\_Size =} \FunctionTok{c}\NormalTok{(n\_before, n\_after)}
\NormalTok{)}

\CommentTok{\# display the table}
\NormalTok{sample\_size}
\end{Highlighting}
\end{Shaded}

\begin{verbatim}
            Stage Sample_Size
1          Before       32561
2 After filtering       32103
\end{verbatim}

The sample size decrease from 32561 to 32103 after restricting. A total
of 458 observations were removed by this restriction.

\subsubsection{e. Remove observations with missing values for
occupation, workclass, or native country. Explain
why.}\label{e.-remove-observations-with-missing-values-for-occupation-workclass-or-native-country.-explain-why.}

\begin{Shaded}
\begin{Highlighting}[]
\CommentTok{\# check the sample size}
\NormalTok{n\_before\_missing }\OtherTok{\textless{}{-}} \FunctionTok{nrow}\NormalTok{(adult)}

\CommentTok{\# remove obs. with missing values in occupation, workclass, or country}
\NormalTok{adult }\OtherTok{\textless{}{-}}\NormalTok{ adult[}
  \FunctionTok{complete.cases}\NormalTok{(adult[, }\FunctionTok{c}\NormalTok{(}\StringTok{"occupation"}\NormalTok{, }\StringTok{"workclass"}\NormalTok{, }\StringTok{"country"}\NormalTok{)]),}
\NormalTok{]}

\CommentTok{\# check the sample size after removing missing values}
\NormalTok{n\_after\_missing }\OtherTok{\textless{}{-}} \FunctionTok{nrow}\NormalTok{(adult)}

\CommentTok{\# display the sample size}
\NormalTok{n\_before\_missing}
\end{Highlighting}
\end{Shaded}

\begin{verbatim}
[1] 32103
\end{verbatim}

\begin{Shaded}
\begin{Highlighting}[]
\NormalTok{n\_after\_missing}
\end{Highlighting}
\end{Shaded}

\begin{verbatim}
[1] 29873
\end{verbatim}

After removing observations with missing values in occupation,
workclass, or country, the sample size decreased from 32,103 to 29,873.
A total of 2,230 observations were removed. Observations with missing
value were removed to ensure complete information for the variables used
in the analysis.

\subsubsection{``https://www.rdocumentation.org/packages/stats/versions/3.6.2/topics/complete.cases''
helped me figure out a more efficient way to code my loop for problem 2
by doing
`complete.cases'.}\label{httpswww.rdocumentation.orgpackagesstatsversions3.6.2topicscomplete.cases-helped-me-figure-out-a-more-efficient-way-to-code-my-loop-for-problem-2-by-doing-complete.cases.}

\subsubsection{f.Eliminate any observations with: age less than 18, age
greater than 90, education level identified as
Preschool.}\label{f.eliminate-any-observations-with-age-less-than-18-age-greater-than-90-education-level-identified-as-preschool.}

\begin{Shaded}
\begin{Highlighting}[]
\CommentTok{\# check the sample size }
\NormalTok{n\_before\_filter }\OtherTok{\textless{}{-}} \FunctionTok{nrow}\NormalTok{(adult)}

\CommentTok{\# remove obs. with age below 18, age above 90, or education level equal to Preschool}
\NormalTok{adult }\OtherTok{\textless{}{-}}\NormalTok{ adult[}
\NormalTok{  adult}\SpecialCharTok{$}\NormalTok{age }\SpecialCharTok{\textgreater{}=} \DecValTok{18} \SpecialCharTok{\&}
\NormalTok{  adult}\SpecialCharTok{$}\NormalTok{age }\SpecialCharTok{\textless{}=} \DecValTok{90} \SpecialCharTok{\&}
\NormalTok{  adult}\SpecialCharTok{$}\NormalTok{education }\SpecialCharTok{!=} \StringTok{"Preschool"}\NormalTok{,}
\NormalTok{]}

\CommentTok{\# check the sample size after filtering}
\NormalTok{n\_after\_filter }\OtherTok{\textless{}{-}} \FunctionTok{nrow}\NormalTok{(adult)}

\CommentTok{\# display}
\NormalTok{n\_before\_filter}
\end{Highlighting}
\end{Shaded}

\begin{verbatim}
[1] 29873
\end{verbatim}

\begin{Shaded}
\begin{Highlighting}[]
\NormalTok{n\_after\_filter}
\end{Highlighting}
\end{Shaded}

\begin{verbatim}
[1] 29523
\end{verbatim}

After removing observations with age less than 18, age greater than 90,
or education level equal to Preschool, the sample size decreased from
29,873 to 29,523.

\subsubsection{g. Report your final sample
size.}\label{g.-report-your-final-sample-size.}

\begin{Shaded}
\begin{Highlighting}[]
\NormalTok{final }\OtherTok{\textless{}{-}} \FunctionTok{nrow}\NormalTok{(adult)}
\NormalTok{final}
\end{Highlighting}
\end{Shaded}

\begin{verbatim}
[1] 29523
\end{verbatim}

The final sample size is 29,523 observations.

\subsection{Q3:Happy Numbers}\label{q3happy-numbers}

\subsubsection{a. Write a function nextHappy() that given a positive
integer, computes the next value in the happy-number sequence. Be sure
to provide a reasonable error on an invalid input. Be sure to document
your
function.}\label{a.-write-a-function-nexthappy-that-given-a-positive-integer-computes-the-next-value-in-the-happy-number-sequence.-be-sure-to-provide-a-reasonable-error-on-an-invalid-input.-be-sure-to-document-your-function.}

\begin{Shaded}
\begin{Highlighting}[]
\CommentTok{\# compute the next value in a happy number sequence}
\NormalTok{nextHappy }\OtherTok{\textless{}{-}} \ControlFlowTok{function}\NormalTok{(n) \{}
  \CommentTok{\# check that n is a single positive integer}
  \ControlFlowTok{if}\NormalTok{(}\FunctionTok{length}\NormalTok{(n) }\SpecialCharTok{!=} \DecValTok{1} \SpecialCharTok{||} \FunctionTok{is.na}\NormalTok{(n) }\SpecialCharTok{||}\NormalTok{ n }\SpecialCharTok{\textless{}=} \DecValTok{0} \SpecialCharTok{||}\NormalTok{ n }\SpecialCharTok{!=} \FunctionTok{floor}\NormalTok{(n)) \{}
    \FunctionTok{stop}\NormalTok{(}\StringTok{"input must be a positive integer"}\NormalTok{)}
\NormalTok{  \}}
  
  \CommentTok{\# split n into individual digits}
\NormalTok{  digits }\OtherTok{\textless{}{-}} \FunctionTok{as.numeric}\NormalTok{(}\FunctionTok{strsplit}\NormalTok{(}\FunctionTok{as.character}\NormalTok{(n), }\StringTok{""}\NormalTok{)[[}\DecValTok{1}\NormalTok{]])}
  
  \CommentTok{\# return the sum of the digits}
  \FunctionTok{sum}\NormalTok{(digits}\SpecialCharTok{\^{}}\DecValTok{2}\NormalTok{)}
\NormalTok{\}}

\CommentTok{\# demonstrate the function using the given examples}
\FunctionTok{nextHappy}\NormalTok{(}\DecValTok{19}\NormalTok{)}
\end{Highlighting}
\end{Shaded}

\begin{verbatim}
[1] 82
\end{verbatim}

\begin{Shaded}
\begin{Highlighting}[]
\FunctionTok{nextHappy}\NormalTok{(}\DecValTok{82}\NormalTok{)}
\end{Highlighting}
\end{Shaded}

\begin{verbatim}
[1] 68
\end{verbatim}

\subsubsection{b.Create a function happySequence() that returns the
sequence beginning at a given number. Use your nextHappy() function to
perform the calculation. Be sure to provide a reasonable error on an
invalid input. Be sure to document your
function.}\label{b.create-a-function-happysequence-that-returns-the-sequence-beginning-at-a-given-number.-use-your-nexthappy-function-to-perform-the-calculation.-be-sure-to-provide-a-reasonable-error-on-an-invalid-input.-be-sure-to-document-your-function.}

\begin{Shaded}
\begin{Highlighting}[]
\CommentTok{\# generate a happy number sequence}
\NormalTok{happySequence }\OtherTok{\textless{}{-}} \ControlFlowTok{function}\NormalTok{(n) \{}
  \CommentTok{\# check that n is a single positive integer}
  \ControlFlowTok{if}\NormalTok{(}\FunctionTok{length}\NormalTok{(n) }\SpecialCharTok{!=} \DecValTok{1} \SpecialCharTok{||} \FunctionTok{is.na}\NormalTok{(n) }\SpecialCharTok{||}\NormalTok{ n }\SpecialCharTok{\textless{}=} \DecValTok{0} \SpecialCharTok{||}\NormalTok{ n }\SpecialCharTok{!=} \FunctionTok{floor}\NormalTok{(n)) \{}
    \FunctionTok{stop}\NormalTok{(}\StringTok{"input must be a positive integer"}\NormalTok{)}
\NormalTok{  \}}
  \CommentTok{\# start the sequence}
\NormalTok{  sequence }\OtherTok{\textless{}{-}}\NormalTok{ n}
  
  \CommentTok{\# continue until the sequence reaches 1 or enters a cycle}
  \ControlFlowTok{while}\NormalTok{ (}\FunctionTok{tail}\NormalTok{(sequence, }\DecValTok{1}\NormalTok{) }\SpecialCharTok{!=} \DecValTok{1}\NormalTok{) \{}
    \CommentTok{\# calculate the next value}
\NormalTok{    next\_value }\OtherTok{\textless{}{-}} \FunctionTok{nextHappy}\NormalTok{(}\FunctionTok{tail}\NormalTok{(sequence, }\DecValTok{1}\NormalTok{))}
    \CommentTok{\#check the next value has already appeared}
    \ControlFlowTok{if}\NormalTok{(next\_value }\SpecialCharTok{\%in\%}\NormalTok{ sequence) \{}
      \FunctionTok{return}\NormalTok{(}\FunctionTok{list}\NormalTok{(}
        \AttributeTok{sequence =}\NormalTok{ sequence,}
        \AttributeTok{isHappy =} \ConstantTok{FALSE}
\NormalTok{      ))}
\NormalTok{    \}}
    \CommentTok{\# add the next value}
\NormalTok{    sequence }\OtherTok{\textless{}{-}} \FunctionTok{c}\NormalTok{(sequence, next\_value)}
\NormalTok{  \}}
  \CommentTok{\# return the sequence and TRUE if reaches 1}
  \FunctionTok{return}\NormalTok{(}\FunctionTok{list}\NormalTok{(}
    \AttributeTok{sequence =}\NormalTok{ sequence,}
    \AttributeTok{isHappy =} \ConstantTok{TRUE}
\NormalTok{  ))}
\NormalTok{\}}

\CommentTok{\# demonstrate the function with a happy number}
\FunctionTok{happySequence}\NormalTok{(}\DecValTok{19}\NormalTok{)}
\end{Highlighting}
\end{Shaded}

\begin{verbatim}
$sequence
[1]  19  82  68 100   1

$isHappy
[1] TRUE
\end{verbatim}

\begin{Shaded}
\begin{Highlighting}[]
\FunctionTok{happySequence}\NormalTok{(}\DecValTok{4}\NormalTok{)}
\end{Highlighting}
\end{Shaded}

\begin{verbatim}
$sequence
[1]   4  16  37  58  89 145  42  20

$isHappy
[1] FALSE
\end{verbatim}

\subsubsection{c.Among integers 1 through 1000: How many are happy? What
percentage are happy? What starting value produces the longest happy
sequence before terminating or
cycling?}\label{c.among-integers-1-through-1000-how-many-are-happy-what-percentage-are-happy-what-starting-value-produces-the-longest-happy-sequence-before-terminating-or-cycling}

\begin{Shaded}
\begin{Highlighting}[]
\CommentTok{\# happy number through 1 to 1000}
\NormalTok{results }\OtherTok{\textless{}{-}} \FunctionTok{lapply}\NormalTok{(}\DecValTok{1}\SpecialCharTok{:}\DecValTok{1000}\NormalTok{, happySequence)}

\CommentTok{\# extract each starting value is happy}
\NormalTok{is\_happy }\OtherTok{\textless{}{-}} \FunctionTok{sapply}\NormalTok{(results, }\ControlFlowTok{function}\NormalTok{(x) x}\SpecialCharTok{$}\NormalTok{isHappy)}

\CommentTok{\# count \# of happy numbers}
\NormalTok{num\_happy }\OtherTok{\textless{}{-}} \FunctionTok{sum}\NormalTok{(is\_happy)}
\NormalTok{num\_happy}
\end{Highlighting}
\end{Shaded}

\begin{verbatim}
[1] 143
\end{verbatim}

\begin{Shaded}
\begin{Highlighting}[]
\CommentTok{\# calculate the percentage}
\NormalTok{per\_happy }\OtherTok{\textless{}{-}} \FunctionTok{mean}\NormalTok{(is\_happy) }\SpecialCharTok{*} \DecValTok{100}
\NormalTok{per\_happy}
\end{Highlighting}
\end{Shaded}

\begin{verbatim}
[1] 14.3
\end{verbatim}

\begin{Shaded}
\begin{Highlighting}[]
\CommentTok{\# calculate the sequence length for each happy number}
\NormalTok{sequence\_lengths }\OtherTok{\textless{}{-}} \FunctionTok{sapply}\NormalTok{(}
\NormalTok{  results,}
  \ControlFlowTok{function}\NormalTok{(x) \{}
    \ControlFlowTok{if}\NormalTok{(x}\SpecialCharTok{$}\NormalTok{isHappy) \{}
      \FunctionTok{length}\NormalTok{(x}\SpecialCharTok{$}\NormalTok{sequence)}
\NormalTok{    \} }\ControlFlowTok{else}\NormalTok{\{}
      \ConstantTok{NA}
\NormalTok{    \}}
\NormalTok{  \}}
\NormalTok{)}

\CommentTok{\# find the max happy sequence length}
\NormalTok{max\_length }\OtherTok{\textless{}{-}} \FunctionTok{max}\NormalTok{(sequence\_lengths, }\AttributeTok{na.rm =} \ConstantTok{TRUE}\NormalTok{)}
\NormalTok{max\_length}
\end{Highlighting}
\end{Shaded}

\begin{verbatim}
[1] 7
\end{verbatim}

\begin{Shaded}
\begin{Highlighting}[]
\CommentTok{\# find the starting values with the longest happy sequence}
\NormalTok{longest\_start }\OtherTok{\textless{}{-}} \FunctionTok{which}\NormalTok{(sequence\_lengths }\SpecialCharTok{==}\NormalTok{ max\_length)}
\NormalTok{longest\_start}
\end{Highlighting}
\end{Shaded}

\begin{verbatim}
 [1] 356 365 379 397 469 496 536 563 635 649 653 694 739 793 899 937 946 964 973
[20] 989 998
\end{verbatim}

Among the integers from 1 to 1000, 143 are happy numbers, which is
14.3\%. The longest happy sequence has length 7. Multiple starting
values produce sequences of this maximum length, including 356, 365,
379, 397, 469, 496, 536, 563, 635, 649, 653, 694, 739, 793, 899, 937,
946, 964, 973, 989, and 998.




\end{document}
